\section*{अध्याय \ref{ch:b2:fourier} --- फूरिये श्रेणियाँ}

\begin{solution}{exo:b2:fourier:1}
विषमता $a_n$ को मार देती है। $n \geq 1$ के लिए:
\[
b_n = \frac{2}{\pi}\int_0^{\pi} \sin nt\,\dd t
= \frac{2}{\pi}\cdot\frac{1 - (-1)^n}{n}
= \begin{cases} \frac{4}{\pi n} & n \text{ विषम},\\ 0 & n
\text{ सम}. \end{cases}
\]
अतः $S(f)(t) = \frac{4}{\pi}\sum_{k\geq0} \frac{\sin\bigl((2k+1)t
\bigr)}{2k+1}$। $t = \frac\pi2$ (कोई संततता-बिंदु, मान $1$) पर डिरिक्ले: $\sin\bigl((2k+1)\frac\pi2\bigr) = (-1)^k$, जिससे
\[
1 = \frac4\pi \sum_{k\geq0}\frac{(-1)^k}{2k+1}
\quad\Longrightarrow\quad
\sum_{k\geq0}\frac{(-1)^k}{2k+1} = \frac{\pi}{4}
\quad\text{(लाइब्नित्स)}.
\]
और उछाल $t = 0$ पर: श्रेणी का योग $0 = \frac{f(0^+) +
f(0^-)}{2}$ है, जैसा डिरिक्ले निर्धारित करती है
(हर पद लुप्त हो जाता है: संगत)।
\end{solution}

\begin{solution}{exo:b2:fourier:2}
समता $b_n$ को मार देती है; $a_0 = \frac{1}{\pi}\int_{-\pi}^\pi\abs
t\,\dd t = \pi$, और $n \geq 1$ के लिए:
\[
a_n = \frac{2}{\pi}\int_0^\pi t\cos nt\,\dd t
= \frac{2}{\pi}\cdot\frac{(-1)^n - 1}{n^2}
= \begin{cases} -\frac{4}{\pi n^2} & n \text{ विषम},\\ 0 & n
\text{ सम}, \end{cases}
\]
(एक खंडशः समाकलन)। अतः
\[
\abs t = \frac{\pi}{2} - \frac{4}{\pi}\sum_{k\geq0}
\frac{\cos\bigl((2k+1)t\bigr)}{(2k+1)^2} ,
\]
और अभिसरण \emph{सामान्य} है ($\sum (2k+1)^{-2} < \infty$) --- जैसा \cref{thm:b2:fourier:parseval}~(1) इस \omterm{def:b2:metric:continuity}{संतत} खंडशः-$C^1$
फलन के लिए बताती है। $t = 0$ पर:
\[
0 = \frac\pi2 - \frac4\pi\sum_{k\geq0}\frac{1}{(2k+1)^2}
\quad\Longrightarrow\quad
\sum_{k\geq0}\frac{1}{(2k+1)^2} = \frac{\pi^2}{8} .
\]
$\sum \frac{1}{n^2}$ को विषम और सम भागों में बाँटने पर: $S =
\frac{\pi^2}{8} + \frac S4$, अतः $S = \frac{\pi^2}{6}$: फिर बाज़ल।
\end{solution}

\begin{solution}{exo:b2:fourier:3}
समता: $b_n = 0$; $a_0 = \frac{1}{\pi}\int_{-\pi}^{\pi} t^2 =
\frac{2\pi^2}{3}$; और दो खंडशः समाकलन $a_n =
\frac{4(-1)^n}{n^2}$ देते हैं ($n \geq 1$)। अतः
\[
t^2 = \frac{\pi^2}{3} + 4\sum_{n\geq1}
\frac{(-1)^n}{n^2}\cos nt
\qquad (\abs t \leq \pi),
\]
जो \omterm{def:b2:funcseq:series}{सामान्य रूप से} अभिसारी है। $t = 0$ पर: $0 = \frac{\pi^2}{3} +
4\sum\frac{(-1)^n}{n^2}$, अर्थात् $\sum_{n\geq1}
\frac{(-1)^{n+1}}{n^2} = \frac{\pi^2}{12}$। पारसेवाल:
\[
\frac{1}{2\pi}\int_{-\pi}^{\pi} t^4\,\dd t = \frac{\pi^4}{5}
= \frac{a_0^2}{4} + \frac12\sum_{n\geq1} a_n^2
= \frac{\pi^4}{9} + 8\sum_{n\geq1}\frac{1}{n^4} ,
\]
अतः $\sum \frac{1}{n^4} = \frac18\bigl(\frac{\pi^4}{5} -
\frac{\pi^4}{9}\bigr) = \frac{\pi^4}{90}$।
\end{solution}

\begin{solution}{exo:b2:fourier:4}
यदि इसके अतिरिक्त $f$ खंडशः $C^1$ हो: तो (सिद्ध) पारसेवाल $\norm f_2^2 = \sum\abs{c_n}^2 = 0$ देती है, और
\omterm{def:b2:metric:continuity}{संतत} $\abs f^2$ के समाकल की यथार्थ धनात्मकता $f = 0$ को बाध्य कर देती है।

और केवल \omterm{def:b2:metric:continuity}{संतत} $f$ के लिए सामान्य पारसेवाल स्वीकार कर लीजिए: वही एक पंक्ति
की उपपत्ति। (सघनता वाले रास्ते की रूपरेखा: फेयेर/वाइरश्ट्रास-प्रकार का
त्रिकोणमितीय सन्निकटन दिखाता है कि \omterm{def:b2:metric:continuity}{संतत} आवर्ती फलनों में त्रिकोणमितीय
बहुपद $\norm\cdot_2$-सघन हैं; और चूँकि $f
\perp$ उन सबको, अतः सन्निकटकों $P$ के लिए $\norm f_2^2 = \langle f, f - P\rangle \leq
\norm f_2\norm{f - P}_2$,
जो $\norm f_2 =
0$ को बाध्य कर देता है।)
\end{solution}

\begin{solution}{exo:b2:fourier:5}
समता: $b_n = 0$;
\[
a_n = \frac{2}{\pi}\int_0^\pi \cos(\alpha t)\cos(nt)\,\dd t
= \frac{2}{\pi}\cdot
\frac{(-1)^n\,\alpha\sin(\pi\alpha)}{\alpha^2 - n^2}
\]
(गुणनफल-से-योग, फिर समाकलन; $a_0 =
\frac{2\sin(\pi\alpha)}{\pi\alpha}$)। $t = \pi$ पर डिरिक्ले (जो आवर्ती विस्तार का
संततता-बिंदु है, क्योंकि समता से उसके एकपक्षीय मान मेल खाते हैं):
\[
\cos(\pi\alpha)
= \frac{\sin(\pi\alpha)}{\pi\alpha}
+ \sum_{n\geq1} \frac{2\alpha\sin(\pi\alpha)}{\pi(\alpha^2 -
n^2)}\,(-1)^n\cos(n\pi)
= \frac{\sin(\pi\alpha)}{\pi}\Bigl(\frac{1}{\alpha} +
\sum_{n\geq1}\frac{2\alpha}{\alpha^2 - n^2}\Bigr),
\]
जिसमें $(-1)^n\cos n\pi = 1$ का उपयोग हुआ। $\sin(\pi\alpha)/\pi$ से भाग देने पर कोस्पर्शज्या का प्रसार मिल
जाता है।
\end{solution}

\begin{solution}{exo:b2:fourier:6}
$c_n(f') = \iu n\,c_n(f)$ को दोहराने पर ($k$ बार, मेल खाते सीमा-मानों के साथ $C^{k}$ टुकड़ों पर
खंडशः समाकलन): $c_n(f^{(k)}) = (\iu n)^k c_n(f)$। और खंडशः \omterm{def:b2:metric:continuity}{संतत} $f^{(k)}$ के गुणांक परिबद्ध हैं (वस्तुतः $\to 0$,
बेसेल):
\[
\abs{c_n(f)} = \frac{\abs{c_n(f^{(k)})}}{\abs n^k}
= O\bigl(\abs n^{-k}\bigr) .
\]
\end{solution}

\begin{solution}{exo:b2:fourier:7}
$f$ तथा $f'$ के लिए पारसेवाल (दोनों वैध: $f$ $C^1$ है, और $f'$ खंडशः \omterm{def:b2:metric:continuity}{संतत} ---
वस्तुतः \omterm{def:b2:metric:continuity}{संतत}):
\[
\frac{1}{2\pi}\int \abs f^2 = \sum_{n\neq0} \abs{c_n}^2
\quad (c_0 = 0 \text{ माध्य-शून्य परिकल्पना से}),
\qquad
\frac{1}{2\pi}\int \abs{f'}^2 = \sum_{n\neq0} n^2\abs{c_n}^2 .
\]
पद-दर-पद $\abs n \geq 1$ के लिए $n^2\abs{c_n}^2 \geq \abs{c_n}^2$: अतः असमिका निकल आती है। और समता सभी $n$ के लिए
$(n^2 - 1)\abs{c_n}^2 = 0$ को बाध्य कर देती है, अर्थात् $\abs n \geq 2$ के लिए $c_n = 0$: $f(t) =
c_1\eu^{\iu t} + c_{-1}\eu^{-\iu t} = a\cos t + b\sin t$ (वास्तविक रूप);
विलोमतः ऐसे $f$ समता देते हैं।
\end{solution}

\begin{solution}{exo:b2:fourier:8}
\cref{exo:b2:fourier:1} से $S_{2N-1}(x) =
\frac4\pi\sum_{k=0}^{N-1}\frac{\sin((2k+1)x)}{2k+1}$। $x_N =
\frac{\pi}{2N}$ पर, $u_k = (2k+1)x_N = \frac{(2k+1)\pi}{2N}$ के साथ:
\[
S_{2N-1}(x_N) = \frac{4}{\pi}\sum_{k=0}^{N-1}\frac{\sin u_k}{2k+1}
= \frac{4}{\pi}\sum_{k=0}^{N-1}\frac{\sin u_k}{u_k}\cdot
\frac{u_k}{2k+1}
= \frac{2}{\pi}\sum_{k=0}^{N-1}\frac{\sin u_k}{u_k}\cdot
\frac{\pi}{N},
\]
क्योंकि $\frac{u_k}{2k+1} = \frac{\pi}{2N}$। बिंदु $u_k$ $\intcc{0}{\pi}$ के $N$ उपअंतरालों (लंबाई $\frac{\pi}{N}$) के मध्यबिंदु
हैं: अतः यह योग \omterm{def:b2:metric:continuity}{संतत} $u \mapsto \frac{\sin u}{u}$ का मध्यबिंदु रीमान योग है, इसलिए वह
\[
\frac{2}{\pi}\int_0^\pi \frac{\sin u}{u}\,\dd u
\approx \frac{2}{\pi}\times 1.8519 \approx 1.179 .
\]
पर अभिसरित होता है। उछाल के पास आंशिक योग मान $1$ से अर्ध-उछाल का $\approx
18\%$
सदा के लिए पार कर जाते हैं: यही गिब्स की परिघटना है, परिमाणित।
\end{solution}

\begin{solution}{exo:b2:fourier:9}
$\cos 3t = 4\cos^3t - 3\cos t$ से:
\[
\cos^3 t = \frac{3\cos t + \cos 3t}{4},
\qquad
\sin^2t\,\cos t = \cos t - \cos^3 t
= \frac{\cos t - \cos 3t}{4} .
\]
हर एक त्रिकोणमितीय बहुपद है, अतः अपनी ही फूरिये श्रेणी के बराबर (गुणांकों
की अद्वितीयता: दो प्रसारों का अंतर ऐसा त्रिकोणमितीय बहुपद होता जिसके सारे
गुणांक शून्य हों)। $\cos^3t$ के लिए: $a_1 = \frac34$, $a_3 =
\frac14$, और शेष सारे $a_n$ तथा सारे $b_n$
शून्य; $c_{\pm1} =
\frac38$, $c_{\pm3} = \frac18$। और $\sin^2t\cos t$ के लिए: $a_1 =
\frac14$, $a_3 = -\frac14$; $c_{\pm1} = \frac18$, $c_{\pm3} =
-\frac18$।
\end{solution}

\begin{solution}{exo:b2:fourier:10}
सीधा परिकलन:
\[
c_n = \frac{1}{2\pi}\int_{-\pi}^{\pi}\eu^{(a - \iu n)t}\dd t
= \frac{\eu^{(a-\iu n)\pi} - \eu^{-(a - \iu n)\pi}}
{2\pi(a - \iu n)}
= \frac{(-1)^n\sinh(a\pi)}{\pi(a - \iu n)} ,
\]
जिसमें $\eu^{\pm\iu n\pi} = (-1)^n$ का उपयोग हुआ। $t = \pi$ पर आवर्ती विस्तार $\eu^{a\pi}$ से $\eu^{-a\pi}$ तक उछलता है;
और डिरिक्ले (सममित आंशिक योग) देती है
\[
\cosh(a\pi) = \sum_{n\in\Z}(-1)^n c_n\,
= \frac{\sinh(a\pi)}{\pi}\Bigl(\frac1a +
\sum_{n\geq1}\Bigl(\frac{1}{a - \iu n} +
\frac{1}{a + \iu n}\Bigr)\Bigr)
= \frac{\sinh(a\pi)}{\pi}\Bigl(\frac1a +
\sum_{n\geq1}\frac{2a}{a^2 + n^2}\Bigr) ,
\]
जहाँ युग्मित पदों के काल्पनिक भाग कट जाते हैं। $\sinh(a\pi)$ से भाग दीजिए:
\[
\coth(\pi a) = \frac{1}{\pi a} +
\sum_{n\geq1}\frac{2a}{\pi(a^2 + n^2)} ,
\]
अर्थात् \cref{exo:b2:fourier:5} का अतिपरवलयिक जुड़वाँ।
\end{solution}

\begin{solution}{exo:b2:fourier:11}
क्रमविनिमेयता: $s = x - t$ प्रतिस्थापित कीजिए और समाकल्य की आवर्तिता का उपयोग
कीजिए। $c_n(f * g)$ के लिए समाकल्य $(x, t) \mapsto
f(x-t)g(t)\eu^{-\iu nx}$ \omterm{def:b2:metric:continuity}{संतत} है; और दोनों क्रमिक समाकल मेल खाते हैं
(बाहरी चर की ऊपरी सीमा के फलनों के रूप में दोनों बाएँ अंत्यबिंदु पर लुप्त
हैं और उनके अवकलज एक ही हैं --- संततता तथा \cref{thm:b2:integration:continuity} क्रमिक समाकल के अवकलन को
न्यायोचित ठहराते हैं)। अतः
\[
c_n(f*g) = \frac{1}{2\pi}\int_{-\pi}^{\pi} g(t)\,\eu^{-\iu nt}
\Bigl(\frac{1}{2\pi}\int_{-\pi}^{\pi}
f(x-t)\,\eu^{-\iu n(x-t)}\dd x\Bigr)\dd t
= c_n(f)\,c_n(g),
\]
जहाँ भीतरी समाकल प्रत्येक $t$ के लिए $c_n(f)$ है (प्रतिस्थापन और आवर्तिता)।
अंत में \cref{lem:b2:fourier:kernel} कहता है $S_N(f)(x) = \frac{1}{2\pi}\int f(x+u)D_N(u)\dd u$; और प्रतिस्थापन $u \mapsto -t$ तथा $D_N$ की समता इसे $(f * D_N)(x)$ में
बदल देते हैं।
\end{solution}

\begin{solution}{exo:b2:fourier:12}
\begin{enumerate}
  \item $w = r\eu^{\iu t}$ के साथ:
        \[
        P_r(t) = 1 + 2\,\Re\frac{w}{1 - w}
        = \Re\frac{1 + w}{1 - w}
        = \frac{1 - \abs w^2}{\abs{1 - w}^2}
        = \frac{1 - r^2}{1 - 2r\cos t + r^2} > 0 .
        \]
        माध्य $1$: सामान्य अभिसारी श्रेणी का पद-दर-पद समाकलन केवल $n = 0$
        बचाता है।
  \item $\delta \leq \abs t \leq \pi$ के लिए: $\cos t \leq
        \cos\delta$, अतः $P_r(t) \leq \frac{1 - r^2}{1 -
        2r\cos\delta + r^2}$, जिसका हर $2 -
        2\cos\delta > 0$ की ओर जाता है जबकि अंश $0$
        की ओर: अर्थात् $0$ के किसी भी प्रतिवेश से बाहर $0$ पर \omterm{def:b2:funcseq:def}{एकसमान
        अभिसरण}।
  \item पद-दर-पद समाकलन ($t$ में \omterm{def:b2:funcseq:series}{सामान्य अभिसरण}) $(f * P_r)(x) = \sum_n r^{\abs
        n}c_n(f)\eu^{\iu nx}$ देता है। सन्निकट
        तत्समक वाला तर्क: माध्य $1$ और धनात्मकता के साथ
        \[
        \abs{(f*P_r)(x) - f(x)}
        \leq \frac{1}{2\pi}\int_{-\pi}^{\pi}
        \abs{f(x-t) - f(x)}\,P_r(t)\,\dd t ,
        \]
        इसे $\abs t = \delta$ पर बाँटिए: अधिकतम $\varepsilon$ (हाइने) जमा $2\norm f_\infty\sup_{\delta\leq\abs
        t\leq\pi}P_r \to \varepsilon$: अतः $r \to 1^-$ होने
        पर \omterm{def:b2:funcseq:def}{एकसमान अभिसरण}। यह फूरिये श्रेणी का \omterm{thm:b2:series:abel}{आबेल योग} है --- घात-श्रेणी
        अध्याय के सीमा-सिद्धांत का फूरिये जुड़वाँ।
\end{enumerate}
\end{solution}

\begin{solution}{pb:b2:fourier:1}
\textbf{1.} $n = 0,
\dots, N-1$ पर \cref{lem:b2:fourier:kernel} का औसत लेने से (समाकल की रैखिकता) $\sigma_N(f)(x) =
\frac{1}{2\pi}\int f(x+u)F_N(u)\dd u$ मिलता है; और
हर $D_n$ का माध्य $1$ है, अतः $F_N$ का माध्य $1$।

\textbf{2.} $\eu^{\iu u} - 1 = 2\iu\,\eu^{\iu
u/2}\sin\frac u2$ के साथ:
\[
\sum_{n=0}^{N-1}\sin\Bigl(\Bigl(n + \frac12\Bigr)u\Bigr)
= \Im\Bigl[\eu^{\iu u/2}\,\frac{\eu^{\iu Nu} - 1}{\eu^{\iu u}
- 1}\Bigr]
= \Re\,\frac{1 - \eu^{\iu Nu}}{2\sin\frac u2}
= \frac{1 - \cos Nu}{2\sin\frac u2}
= \frac{\sin^2\frac{Nu}2}{\sin\frac u2} .
\]
$N\sin\frac u2$ से भाग देने पर:
\[
F_N(u) = \frac1N\sum_{n=0}^{N-1}
\frac{\sin\bigl((n+\frac12)u\bigr)}{\sin\frac u2}
= \frac{1}{N}\,
\frac{\sin^2\frac{Nu}{2}}{\sin^2\frac u2} \geq 0 .
\]

\textbf{3.} $\delta \leq \abs u \leq \pi$ पर: $\sin^2\frac u2
\geq \sin^2\frac\delta2$ और $\sin^2\frac{Nu}2 \leq 1$: अतः वहाँ \omterm{def:b2:funcseq:def}{एकसमान रूप से} $F_N
\leq \frac{1}{N\sin^2(\delta/2)} \to 0$।

\textbf{4.} इकाई माध्य से $\sigma_Nf(x) - f(x) =
\frac{1}{2\pi}\int\bigl(f(x+u) - f(x)\bigr)F_N(u)\dd u$। $\varepsilon$ दिया हो तो हाइने $\abs{f(x+u) - f(x)}
\leq \varepsilon$ वाला $\delta$ देती है
($\abs u \leq \delta$ के लिए, $x$ में एकसमान)। फिर $F_N \geq 0$ तथा उसके इकाई माध्य का उपयोग करके
बड़े $N$ के लिए, $x$ में \omterm{def:b2:funcseq:def}{एकसमान रूप से},
\[
\abs{\sigma_Nf(x) - f(x)}
\leq \varepsilon +
2\norm f_\infty\cdot\frac{1}{2\pi}
\int_{\delta\leq\abs u\leq\pi}F_N
\leq \varepsilon + \frac{2\norm
f_\infty}{N\sin^2\frac\delta2}
\leq 2\varepsilon
\]
अर्थात् फेयेर की प्रमेय।

\textbf{5.} $F_N$ सम है और उसका माध्य $1$ है: अतः हर आधा $\intcc{0}{\pi}$, $\intcc{-\pi}{0}$ माध्य $\frac12$
वहन करता है। तब
\[
\sigma_Nf(x) - \frac{f(x^+)+f(x^-)}{2}
= \frac{1}{2\pi}\int_0^\pi\bigl(f(x+u) -
f(x^+)\bigr)F_N\,\dd u
+ \frac{1}{2\pi}\int_{-\pi}^0\bigl(f(x+u) -
f(x^-)\bigr)F_N\,\dd u ;
\]
हर आधे पर, $\abs u = \delta$ पर बाँटिए जहाँ एकपक्षीय सीमा $\varepsilon$-निकट है, और दूर वाले भाग
को प्रश्न 3 से मार दीजिए: दोनों समाकल $0$ की ओर जाते हैं। परिबंध: $F_N \geq 0$ $\abs{\sigma_Nf(x)} \leq
\frac{1}{2\pi}\int\abs{f(x+u)}F_N \leq \norm f_\infty$
देता है।

\textbf{6.} हर $\sigma_N(f)$ त्रिकोणमितीय बहुपद है ($S_n(f)$, $n < N$ का औसत), और \omterm{def:b2:funcseq:def}{एकसमान रूप
से} $\sigma_N(f) \to f$: अर्थात् त्रिकोणमितीय वाइरश्ट्रास प्रमेय।

\textbf{7.} सभी $n$ के लिए $c_n(f) = 0$ हर $S_n(f) = 0$ को, अतः हर $\sigma_N(f) = 0$ को शून्य बना देता है;
और फेयेर से $f = \lim\sigma_Nf =
0$। इसे किसी अंतर पर लगाइए: \omterm{def:b2:metric:continuity}{संतत} आवर्ती फलन अपने फूरिये
गुणांकों से निर्धारित हो जाते हैं।

\textbf{8.} $S_Nf$ $f$ का $\mathcal T_N$ (\cref{prop:b2:fourier:bessel}) पर लांबिक प्रक्षेप है, अतः वह $P \in \mathcal T_N$ पर $\norm{f - P}_2$
को न्यूनतम करता है; और चूँकि $\sigma_Nf
\in \mathcal T_{N-1} \subseteq \mathcal T_N$:
\[
\norm{f - S_Nf}_2 \leq \norm{f - \sigma_Nf}_2
\leq \norm{f - \sigma_Nf}_\infty \to 0
\]
(बीच वाली असमिका इसलिए कि $\abs\cdot^2$ का माध्य अधिकतम उच्चतम का वर्ग है)।
पाइथागोरस $\norm f_2^2 = \sum_{\abs
n\leq N}\abs{c_n}^2 + \norm{f - S_Nf}_2^2$ फिर सीमा तक चला जाता है: अर्थात् प्रत्येक \omterm{def:b2:metric:continuity}{संतत} आवर्ती $f$
के लिए पारसेवाल।

\textbf{9.} मान लीजिए $t_1, \dots, t_p$ किसी आवर्तकाल में $f$ के उछाल हैं, $M = \norm f_\infty$। छोटे $\eta$
के लिए $g = f$ को अंतरालों $\intoo{t_j - \eta}{t_j + \eta}$ के बाहर वही रखिए और हर ऐसे अंतराल पर आनत जीवा
से परिभाषित कीजिए: तब $g$ \omterm{def:b2:metric:continuity}{संतत}, आवर्ती है, $\norm g_\infty \leq M$, और छोटे $\eta$ के लिए
\[
\norm{f - g}_2^2 \leq \frac{1}{2\pi}\,p\cdot(2M)^2\cdot2\eta
\leq \varepsilon^2
\]
बेसेल $S_N$ को $\norm\cdot_2$ के लिए संकुचन बना देता है, अतः
\[
\norm{f - S_Nf}_2
\leq \norm{f - g}_2 + \norm{g - S_Ng}_2 + \norm{S_N(g -
f)}_2
\leq 2\varepsilon + \norm{g - S_Ng}_2 ,
\]
और प्रश्न 8 प्रत्येक $\varepsilon$ के लिए $\limsup_N\norm{f - S_Nf}_2 \leq
2\varepsilon$ देता है: $\norm{f - S_Nf}_2 \to
0$, और पाइथागोरस प्रत्येक
खंडशः \omterm{def:b2:metric:continuity}{संतत} $f$ के लिए पारसेवाल दे देता है: इस प्रकार \cref{thm:b2:fourier:parseval} का ``स्वीकृत''
अब प्रमेय है।

\textbf{10.} वर्ग तरंग के लिए $\norm f_\infty = 1$, अतः प्रश्न 5 प्रत्येक $N$ के लिए और
सर्वत्र $\abs{\sigma_N(f)} \leq 1$ देता है --- कोई अतिक्रमण कभी नहीं --- जबकि \cref{exo:b2:fourier:8} दिखाता है कि
$\sup_xS_{2N-1}(f)(x) \to
\approx 1.179$। उत्तरदायी गुणधर्म: $F_N \geq 0$, अतः $\sigma_Nf(x)$ $f$ के मानों का भारित \emph{औसत} है
और $\intcc{\min f}{\max f}$ से कभी बाहर नहीं जा सकता; जबकि $D_N$ ऋणात्मक मान लेता है, इसलिए $S_N$
जा सकता है।

\textbf{11.} आगमन से $\abs{\sin N\theta} \leq N\abs{\sin\theta}$ ($\abs{\sin(N{+}1)\theta} \leq
\abs{\sin N\theta}\abs{\cos\theta} +
\abs{\cos N\theta}\abs{\sin\theta} \leq
(N+1)\abs{\sin\theta}$): $\theta = \frac u2$ के साथ,
\[
F_N(u) = \frac{\sin^2\frac{Nu}2}{N\sin^2\frac u2}
\leq \frac{N^2\sin^2\frac u2}{N\sin^2\frac u2} = N .
\]
और $\intcc{0}{\frac\pi2}$ पर $\sin$ की अवतलता $0 \leq u \leq \pi$ के लिए $\sin\frac u2 \geq \frac{u}{\pi}$ देती है, अतः $F_N(u) \leq \frac{1}{N(u/\pi)^2} = \frac{\pi^2}{Nu^2}$।

\textbf{12.} समता तथा $\frac1N$ पर विभाजन से:
\[
\frac{1}{2\pi}\int_{-\pi}^{\pi}\abs uF_N
= \frac1\pi\int_0^\pi uF_N
\leq \frac1\pi\Bigl(\int_0^{1/N}uN\,\dd u +
\int_{1/N}^{\pi}\frac{\pi^2}{Nu}\,\dd u\Bigr)
= \frac{1}{2\pi N} + \frac{\pi\ln(\pi N)}{N} .
\]
$N \geq 2$ के लिए: $\ln(\pi N) \leq \bigl(1 +
\frac{\ln\pi}{\ln2}\bigr)\ln N \leq 2.66\ln N$ और $\frac{1}{2\pi N} \leq \frac{\ln N}{N}$, अतः आघूर्ण $\leq \frac{9\ln N}{N}$ है।

\textbf{13.} $L$-लिप्शिट्स $f$ के लिए:
\[
\abs{\sigma_Nf(x) - f(x)}
\leq \frac{1}{2\pi}\int\abs{f(x+u) - f(x)}F_N(u)\dd u
\leq L\cdot\frac{1}{2\pi}\int\abs uF_N
\leq \frac{9L\ln N}{N},
\]
और यह $x$ में एकसमान है।

\textbf{14.} $c_0(e_1) = 0$ $S_0(e_1) = 0$ देता है, जबकि $n \geq 1$ के लिए $S_n(e_1) = e_1$: $\sigma_N(e_1) =
\frac{N-1}{N}e_1$, अतः $\norm{\sigma_Ne_1 - e_1}_\infty =
\frac1N$। इस
\omterm{def:b2:metric:complete}{पूर्ण}, बैंड-सीमित संकेत के लिए भी दर $\frac1N$ है: अर्थात् फेयेर संतृप्त हो
जाता है, ठीक वैसे ही जैसे बर्नस्टाइन संकारक $\frac1n$ पर संतृप्त होता है
(वोरोनोव्स्काया, फलन-अनुक्रम अध्याय की सप्ताहांत समस्या में)।

\textbf{15.} यदि $f$ $\intoo{x-\delta}{x+\delta}$ पर लुप्त हो, तो $\sigma_Nf(x) = \frac{1}{2\pi}\int_{\delta \leq \abs u
\leq \pi}f(x+u)F_N(u)\dd u$, जिसका निरपेक्ष मान अधिकतम
$\frac{\norm f_\infty}{N\sin^2(\delta/2)} \to 0$ है: अतः $x$ पर चेज़ारो माध्य केवल $x$ के पास वाले $f$ को देखते हैं।

\textbf{16.} $\intcc ab$ और $\varepsilon$ दिए हों तो ऐसे \omterm{def:b2:metric:continuity}{संतत} $1$-आवर्ती खंडशः-आनत $\varphi^\pm$
चुनिए कि $\varphi^- \leq \mathbf 1_{\intcc ab} \leq \varphi^+$ और $\int_0^1(\varphi^+ - \varphi^-) \leq \varepsilon$ हो (कुल लंबाई $\varepsilon$ वाले अंतरालों पर ढाल वाले समलंब)।
तब
\[
\limsup_N\frac{\#\{n \leq N : x_n \in \intcc ab\}}{N}
\leq \lim_N\frac1N\sum_{n\leq N}\varphi^+(x_n)
= \int_0^1\varphi^+ \leq b - a + \varepsilon ,
\]
और सममित रूप से $\liminf \geq b - a - \varepsilon$: अतः अनुपात $b - a$ की ओर जाता है।

\textbf{17.} $k \neq 0$ वाले $f = \eu^{2\iu\pi k\cdot}$ के लिए परिकल्पना सीमा $0 = \int_0^1f$ देती है; और $k = 0$ के
लिए दोनों पक्ष $1$ हैं; तथा रैखिकता हर त्रिकोणमितीय बहुपद निपटा देती है।
\omterm{def:b2:metric:continuity}{संतत} $1$-आवर्ती $f$ और $\varepsilon >
0$ के लिए प्रश्न 6 ($t = 2\pi x$ से ले जाकर) $\norm{f - P}_\infty \leq
\varepsilon$ वाला कोई
त्रिकोणमितीय बहुपद $P$ दे देता है:
\[
\Bigl|\frac1N\sum_{n\leq N}f(x_n) - \int_0^1f\Bigr|
\leq 2\varepsilon +
\Bigl|\frac1N\sum_{n\leq N}P(x_n) - \int_0^1P\Bigr|
\longrightarrow 2\varepsilon .
\]
प्रश्न 16 के साथ: $(x_n)$ समवितरित है।

\textbf{18.} $k \neq 0$ और अपरिमेय $\alpha$ के लिए $w =
\eu^{2\iu\pi k\alpha} \neq 1$:
\[
\Bigl|\sum_{n=1}^{N}w^n\Bigr|
= \Bigl|\frac{w(w^N - 1)}{w - 1}\Bigr|
\leq \frac{2}{\abs{1 - w}} ,
\]
और यह परिबंध $N$ से स्वतंत्र है; $N$ से भाग देने पर वाइल की कसौटी मिल
जाती है, और प्रश्न 17 निष्कर्ष दे देता है: $(\{n\alpha\})$ समवितरित है।

\textbf{19.} हर उपअंतराल को उसकी लंबाई के बराबर अनंतस्पर्शी अनुपात मिलता
है, विशेष रूप से अपरिमित बिंदु: अतः $(\{n\alpha\})$ सघन है। समवितरण अधिक मज़बूत है:
कोई अनुक्रम सघन होते हुए भी अपना लगभग सारा समय किसी एक कोने में बिता सकता
है (सघनता कहती है कि अनुक्रम \emph{कहाँ} जाता है, समवितरण कहता है
\emph{कितनी बार})।

\textbf{20.} $2^n$ का अग्र अंक $1$ है यदि और केवल यदि किसी $m$ के लिए $10^m \leq 2^n <
2\cdot10^m$,
अर्थात् यदि और केवल यदि $\{n\log_{10}2\} \in
\intco{0}{\log_{10}2}$। अपरिमेयता: $\log_{10}2 =
\frac pq$ से $2^q = 10^p = 2^p5^p$ मिलता, जो अनन्य
गुणनखंडन के कारण $p
\geq 1$ के लिए असंभव है। वाइल की प्रमेय (प्रश्न 18, $\alpha = \log_{10}2$ के
साथ; अर्ध-विवृत अंतराल को समीप की लंबाइयों वाले संवृत अंतरालों के बीच
निचोड़ा जाता है) अनुपात $\log_{10}2 \approx 0.301$ दे देती है: अर्थात् $2^n$ के पहले अंक बेनफर्ड
के नियम का पालन करते हैं।

\textbf{21.} $z$ $C^1$ है, अतः उसकी फूरिये श्रेणी \omterm{def:b2:funcseq:series}{सामान्य रूप से} अभिसरित
होती है और उसका योग $z$ है (\cref{thm:b2:fourier:parseval}~(1)), तथा $c_n(z') = \iu n\,c_n$। और चूँकि $\abs{z'} = \frac{L}{2\pi}$ अचर है,
\[
\int_0^{2\pi}\abs{z'}^2\dd t
= 2\pi\Bigl(\frac{L}{2\pi}\Bigr)^{\!2}
= \frac{L^2}{2\pi} ,
\]
और \omterm{def:b2:metric:continuity}{संतत} $z'$ पर लगाई गई पारसेवाल $\frac{1}{2\pi}\int_0^{2\pi}\abs{z'}^2 =
\sum_n\abs{\iu nc_n}^2$ देती है, अर्थात् $\frac{L^2}{2\pi} = 2\pi\sum_n
n^2\abs{c_n}^2$।

\textbf{22.} ध्रुवित पारसेवाल $\frac{1}{2\pi}\int\conj
fg = \sum\conj{c_n(f)}c_n(g)$ $f + g$ और $f + \iu g$ पर लगाई गई पारसेवाल से निकलती
है (\omterm{def:b2:quadratic:def}{ध्रुवण सर्वसमिका}), और दोनों \omterm{def:b2:metric:continuity}{संतत} हैं। $f = z$, $g = z'$ के साथ:
\[
A = \frac12\,\Im\int_0^{2\pi}\conj z\,z'
= \pi\,\Im\sum_n\conj{c_n}(\iu n c_n)
= \pi\sum_n n\abs{c_n}^2 .
\]

\textbf{23.} प्रश्न 21--22 मिलाने पर:
\[
L^2 - 4\pi A
= 4\pi^2\sum_n n^2\abs{c_n}^2 -
4\pi^2\sum_n n\abs{c_n}^2
= 4\pi^2\sum_{n\in\Z}(n^2 - n)\abs{c_n}^2 \geq 0 ,
\]
क्योंकि हर पूर्णांक के लिए $n^2 - n = n(n-1) \geq 0$। और समता सभी $n \notin \{0, 1\}$ के लिए $c_n = 0$ को बाध्य कर
देती है: $z(t) = c_0 +
c_1\eu^{\iu t}$, अर्थात् केंद्र $c_0$ और त्रिज्या $\abs{c_1} = \frac{L}{2\pi}$ वाला वृत्त (अचर चाल)।
यही समपरिमापी असमिका की हुर्विट्स वाली उपपत्ति है: $A \leq \frac{L^2}{4\pi}$, और केवल वृत्त।

\textbf{24.} त्रिज्या $R$ का वृत्त: $L = 2\pi R$, $A = \pi
R^2$: $L^2 = 4\pi^2R^2 = 4\pi A$: अर्थात् समता। भुजा $a$
का वर्ग: $L^2 = 16a^2 > 4\pi a^2 = 4\pi A$ ($16 > 4\pi \approx
12.57$ के रूप में)। और ऋणात्मक $n$ के लिए $n^2 - n = n(n - 1)$ दो ऋणात्मक
पूर्णांकों का गुणनफल है: अर्थात् धनात्मक --- इसलिए पीछे घूमने वाली विधाएँ
क्षेत्रफल की दुगुनी क़ीमत लेती हैं। अचर चाल प्रश्न 21 में आई, जहाँ उसने
$\int\abs{z'}^2$ को $\frac{L^2}{2\pi}$ में बदला; और अनचर चाल के लिए कोशी--श्वार्ज़ $\int\abs{z'}^2 \geq
\frac{(\int\abs{z'})^2}{2\pi} = \frac{L^2}{2\pi}$ देता है, अतः
असमिका बची रहती है और समता की एकमात्र स्थिति अब भी वृत्त ही है।

\textbf{25.} (क) सब कुछ $F_N \geq 0$ से बहता है (इकाई माध्य और संकेंद्रण के साथ);
$D_N$ का माध्य इकाई है और दोलन का संकेंद्रण भी, पर धनात्मकता नहीं, और गिब्स
उसकी क़ीमत है। (ख) फूरिये श्रेणी की चेज़ारो योग्यता से उसी योग के साथ उसकी
आबेल योग्यता निकलती है (फ्रोबेनियस, जो घात-श्रेणी अध्याय की सप्ताहांत
समस्या में सिद्ध है) --- और \cref{exo:b2:fourier:12} का प्वासों-कर्नेल वाला रास्ता ठीक आबेल की
विधि ही है। (ग) वाइल की प्रमेय को केवल एकसमान सन्निकटन चाहिए था (प्रश्न
6); जबकि समपरिमापी असमिका को स्वयं पारसेवाल चाहिए थी (प्रश्न 8, 21--22)।
(घ) तृतीय वर्ष का खंड पूर्णता सिद्ध करता है: चरघातांकी $L^2$ का हिल्बर्ट
आधार बनाते हैं, पारसेवाल हिल्बर्ट समष्टियों की समदूरिकता बन जाती है, और
फेयेर की प्रमेय यह कथन बन जाती है कि वह समदूरिकता धनात्मक औसतों से परिकलनीय
है।
\end{solution}
